% Source: https://artofproblemsolving.com/community/c3h2858991_work_done
% Stable identity: aops-topic-2858991-q01
Together they do $\dfrac1p$ of the job per day, so after $q$ days $\dfrac qp$ is done and $\dfrac{p-q}p$ remains. B does this in $r$ days, so B's rate is $\dfrac{p-q}{pr}$ and B alone takes $\boxed{\dfrac{pr}{p-q}}$ days.

A's rate is $\dfrac1p-\dfrac{p-q}{pr}=\dfrac{r-p+q}{pr}$, so A alone takes $\boxed{\dfrac{pr}{q+r-p}}$ days (this needs $q+r>p$). A works only during the first $q$ days, doing $\boxed{\dfrac{q(q+r-p)}{pr}}$ of the job; B does the rest.
